% 半人马小行星
% license CCBYSA3
% type Wiki

本文根据 CC-BY-SA 协议转载翻译自维基百科\href{https://en.wikipedia.org/wiki/Centaur_(small_Solar_System_body)}{相关文章}。

\begin{figure}[ht]
\centering
\includegraphics[width=6cm]{./figures/4207f1463529dc38.png}
\caption{} \label{fig_Centau_2}
\end{figure}
在行星天文学中，半人马天体（centaur）是指一类在木星与海王星之间绕太阳运行、并且其轨道会与一颗或多颗巨行星轨道相交的小型太阳系天体。由于这一特性，半人马天体的轨道通常是不稳定的；几乎所有半人马天体的轨道动力学寿命都只有数百万年$^{[1]}$，但有一个已知的例外，即 514107 Kaʻepaokaʻāwela，其轨道可能是稳定的（尽管是逆行轨道）$^{[2]}$（见注释$^{[1]}$）。半人马天体通常同时表现出小行星和彗星的特征。它们的名称来源于神话中的半人马，即兼具马与人形态的生物。由于观测上对大尺寸天体存在偏向性，半人马天体的总体数量难以确定。对直径大于 1 km 的半人马天体数量估计范围从约 44,000 个$^{[1]}$到超过 1,000 万个$^{[4][5]}$不等。

按照喷气推进实验室（JPL）的定义（也是本文采用的定义），第一个被发现的半人马天体是 944 Hidalgo，发现于 1920 年。然而，直到 1977 年发现 2060 Chiron 之后，半人马天体才被视为一个独立的天体族群。目前已确认的最大半人马天体是 10199 Chariklo，其直径约为 250 km，大小相当于一颗中等尺寸的主带小行星，并且已知其拥有一套环系统。该天体发现于 1997 年。

目前尚无任何半人马天体被近距离拍摄过，不过有证据表明，土星的卫星菲比（Phoebe）——由卡西尼号探测器于 2004 年拍摄——可能是一颗起源于柯伊伯带并被俘获的半人马天体$^{[6]}$。此外，哈勃空间望远镜也获取了一些关于 8405 Asbolus 表面特征的信息。

谷神星可能起源于外行星区域$^{[7]}$，如果这一假设成立，它可以被视为一颗前半人马天体，但目前观测到的半人马天体均起源于其他区域。

在已知占据半人马式轨道的天体中，约有 30 个显示出类似彗星的尘埃彗发，其中 2060 Chiron、60558 Echeclus 和 29P/Schwassmann–Wachmann 1 在完全位于木星轨道之外的情况下仍表现出可探测的挥发物释放$^{[8]}$。因此，Chiron 和 Echeclus 同时被归类为半人马天体和彗星，而 Schwassmann–Wachmann 1 一直被归类为彗星。其他半人马天体（如 52872 Okyrhoe）也被怀疑曾出现过彗发。理论上，任何被扰动并足够接近太阳的半人马天体，最终都可能演化为彗星。
\subsection{分类}
若一颗天体的近日点距离或半长轴位于外行星之间（即木星与海王星之间），即可被归类为半人马天体。由于该区域内轨道在长期尺度上的固有不稳定性，即便是当前并未与任何行星轨道相交的半人马天体（如 2000 GM137 和 2001 XZ255），其轨道也会逐渐演化，并最终受到扰动，从而开始与一颗或多颗巨行星的轨道相交$^{[1]}$。一些天文学家仅将半长轴位于外行星区域的天体归为半人马天体；而另一些学者则认为，只要近日点位于该区域内即可归类为半人马天体，因为这两类轨道在动力学上同样是不稳定的。
\subsubsection{判据差异}
然而，不同机构在对边界天体进行分类时采用了不同的判据，这些判据通常基于轨道要素的具体取值：
\begin{itemize}
\item 小行星中心（Minor Planet Center，MPC）将半人马天体定义为近日点位于木星轨道之外（5.2 AU < q），且半长轴小于海王星半长轴（a < 30.1 AU）的天体$^{[9]}$。MPC 有时会将半人马天体与散射盘天体列为同一类群$^{[10]}$。
\item 喷气推进实验室（Jet Propulsion Laboratory，JPL）采用了相近的定义，将半人马天体界定为半长轴 a 介于木星与海王星之间的天体（5.5 AU ≤ a ≤ 30.1 AU）$^{[11]}$。
\item 相比之下，深黄道巡天（Deep Ecliptic Survey，DES）使用的是一种动力学分类方案。该分类基于将当前轨道向未来延拓 1000 万年时，其轨道行为在数值模拟中的变化。DES 将半人马天体定义为非共振天体，其瞬时（密切）近日点在模拟过程中任意时刻小于海王星的瞬时半长轴。该定义旨在与“行星交叉轨道”等价，并暗示这类天体在当前轨道上的寿命相对较短$^{[12]}$。
\item 《海王星之外的太阳系》（The Solar System Beyond Neptune，2008）一书将半长轴位于木星与海王星之间、且相对于木星的 Tisserand 参数大于 3.05 的天体定义为半人马天体；将相对于木星的 Tisserand 参数低于该阈值、并且为排除柯伊伯带天体而额外施加一个任意的近日点限制（q ≤ 7.35 AU，位于通往土星距离的一半处）的天体归类为木星族彗星；而将半长轴大于海王星、且处于不稳定轨道上的天体归类为散射盘天体$^{[13]}$。
\item 还有一些天文学家倾向于将半人马天体定义为：不与行星处于共振状态、近日点位于海王星轨道以内、并且可以证明在未来 1000 万年内很可能进入某一气体巨行星希尔球范围的天体$^{[14]}$。在这种定义下，半人马天体可被视为向内散射的天体，其与行星的相互作用更强、轨道演化也比典型的散射盘天体更快。
\item JPL 小天体数据库目前列出了 910 颗半人马天体$^{[15]}$。此外，还有 223 个跨海王星天体（即半长轴大于海王星，30.1 AU ≤ a），其近日点距离小于天王星轨道（q ≤ 19.2 AU）$^{[16]}$。
\end{itemize}
\subsubsection{歧义天体}
Gladman 与 Marsden（2008）$^{[13]}$ 的判据会将一些天体归类为木星族彗星：Echeclus（q = 5.8 AU，$T_J = 3.03$）和 Okyrhoe（q = 5.8 AU，$T_J = 2.95$）在传统上一直被归类为半人马天体。Hidalgo（q = 1.95 AU，$T_J = 2.07$）传统上被视为一颗小行星，但在 JPL 的分类中被列为半人马天体，在该判据下也会改变类别而被归为木星族彗星。29P/Schwassmann–Wachmann（q = 5.72 AU，$T_J = 2.99$）则根据所采用的定义不同，被同时归类为半人马天体或木星族彗星$^{[]}$。

其他处于不同分类方法分歧中的天体还包括 (44594) 1999 OX3，其半长轴为 32 AU，但会穿越天王星和海王星的轨道。在深黄道巡天（DES）中，它被列为外侧半人马天体。在内侧半人马天体中，(434620) 2005 VD 的近日点距离非常接近木星，其在 JPL 和 DES 中均被列为半人马天体。

一项近期关于柯伊伯带天体经由半人马区域演化的轨道模拟研究$^{[4]}$ 识别出一个位于 5.4–7.8 AU 之间、寿命较短的“轨道通道”，约有 21\% 的半人马天体会经过该区域，其中包括 72\% 最终演化为木星族彗星的半人马天体。目前已知有 4 个天体位于这一通道区域内（29P/Schwassmann–Wachmann、P/2010 TO20 LINEAR-Grauer、P/2008 CL94 Lemmon 以及 2016 LN8），但数值模拟表明，可能还存在数量级约为 1000 个、半径大于 1 km、尚未被探测到的天体。处于该通道区域内的天体可能表现出显著的活动性$^{[17][18]}$，并且正处于一个重要的演化过渡状态，这进一步模糊了半人马天体与木星族彗星两类天体之间的界限。
\subsection{命名惯例}
根据国际天文学联合会小天体命名工作组（Working Group for Small Bodies Nomenclature，WGSBN；原小天体命名委员会）$^{[19]}$ 的规定，半长轴小于 30 AU（位于海王星轨道以内）且近日点距离大于 5.5 AU（位于木星轨道之外）的半人马天体，应以神话中的半人马来命名，即半人半马的神话生物$^{[20]}$。WGSBN 在技术上将穿越海王星轨道的跨海王星天体（半长轴大于 30 AU、近日点小于 30 AU）也计入半人马天体的范畴，但为其保留了另一套不同的命名方案$^{[20]}$。该方案于 2007 年被正式采用，当时首批此类天体——65489 Ceto–Phorcys 和 42355 Typhon–Echidna——被命名$^{[19]}$。

根据 WGSBN 的规定，穿越海王星轨道的跨海王星天体必须以神话中的奇美拉类生物命名$^{[20]}$，这类生物包括具有混合形态或可变形态的神话生物$^{[19]}$。这一类别中的一个命名实例是 471325 Taowu，其名称源自中国神话中的“饕餮”，据说是一种兼具人、虎和野猪特征的混合生物$^{[21]}$。
\subsection{轨道}
\subsubsection{分布}
\begin{figure}[ht]
\centering
\includegraphics[width=6cm]{./figures/acc469512e466fac.png}
\caption{已知半人马小行星的轨道分布$^{[note 2]}$} \label{fig_Centau_3}
\end{figure}
该示意图展示了已知半人马天体的轨道相对于行星轨道的位置关系。对于部分选定的天体，其轨道偏心率以红色线段表示（从近日点延伸至远日点）。

半人马天体的轨道偏心率分布范围很广，从高度偏心的轨道（如 Pholus、Asbolus、Amycus、Nessus）到接近圆形的轨道（如 Chariklo 以及穿越土星轨道的 Thereus 和 Okyrhoe）。

为说明轨道参数的多样性，图中以黄色标出了若干具有极端轨道特征的天体：
\begin{itemize}
\item 1999 XS35（阿波罗型小行星）具有极端的高偏心率轨道（(e = 0.947)），其轨道范围从地球轨道以内（0.94 AU）一直延伸到远超海王星轨道之外（$>34 AU$）。
\item 2007 TB434 具有近乎圆形的轨道（$(e < 0.026)$）。
\item 2001 XZ255 具有最低的轨道倾角（$(i < 3^\circ)$）。
\item 2004 YH32 是少数具有极端顺行高倾角（$(i > 60^\circ)$）的半人马天体之一，其轨道倾角高达 $(79^\circ)$。尽管其轨道在太阳距离上从小行星带附近延伸至土星轨道以外，但若将其轨道投影到木星轨道平面上，甚至不会到达木星的距离。
\end{itemize}
已有十余个已知半人马天体沿逆行轨道运行，其轨道倾角从中等值（例如 Dioretsa 的约 $(160^\circ)$）,到极端值（$(i>120^\circ)$，例如 (342842) 2008 YB3 的约,$(105^\circ)$ $^{[22]}$）。其中有十七个高倾角逆行半人马天体曾被有争议地认为可能具有星际起源 $^{[23][24][25]}$。
\subsubsection{轨道变化}
\begin{figure}[ht]
\centering
\includegraphics[width=8cm]{./figures/8ef3021d0bf5acc7.png}
\caption{在未来 5500 年内，Asbolus 的半长轴变化情况，采用了两组略有不同的当前轨道要素估计进行计算。在公元 4713 年与木星发生近距离遭遇之后，这两种计算结果开始出现分歧。$^{[26]}$} \label{fig_Centau_4}
\end{figure}
由于半人马天体不受轨道共振的保护，它们的轨道在 $10^{6}-10^{7}$ 年的时间尺度内是不稳定的。$^{[27]}$ 例如，55576 Amycus 位于接近天王星 3:4 共振的位置，其轨道是不稳定的。$^{[1]}$ 对其轨道的动力学研究表明，处于“半人马天体”状态很可能是天体从柯伊伯带向木星族短周期彗星过渡过程中的一个中间轨道阶段。小行星 ((679997)) 2023 RB 将在 2201 年与土星的一次近距离接近中，其轨道发生显著改变。

天体可能首先从柯伊伯带受到扰动，从而成为穿越海王星轨道的天体，并与该行星发生引力相互作用（见其起源理论）。随后它们被归类为半人马天体，但其轨道呈现混沌特征，在与一个或多个外行星反复发生近距离接近的过程中迅速演化。一些半人马天体会进一步演化为穿越木星轨道的天体，其近日点可能被降低至内太阳系；若此时它们表现出彗星活动，则可能被重新归类为木星族活动彗星。最终，半人马天体要么与太阳或某一行星发生碰撞，要么在与行星（尤其是木星）的一次近距离接近后被抛射进入星际空间。
\subsection{物理特性}
\begin{figure}[ht]
\centering
\includegraphics[width=6cm]{./figures/7822bdb4c4d1ac8a.png}
\caption{土星的卫星菲比（Phoebe）于 2004 年由卡西尼–惠更斯号探测器成像，理论上被认为是一颗被俘获的半人马天体。} \label{fig_Centau_5}
\end{figure}